\chapter {Introduction}

Logistics is one of the most important functions of a company
since it affects all other activities in the company. 
%Every system needs a simple or complicated, physical or non-physical 
%distribution network on which goods, services and information move 
%between specified origin and destination points to satisfy customer demands. 
%Logistics designs the distribution network, 
Logistics plans and controls the efficiency of the flow on the network and 
the efficiency of the storage activites in the system.
Logistics costs include three main cost components: warehousing costs including 
inventory carrying costs, transportation costs and logistics administration costs. 
%-economic importance cannot be overstated
Based on estimates for the U.S. in 2002, transportation costs were \$577 billion, 
inventory carrying costs and warehousing costs were \$298 billion 
and logistics administration costs were \$35 billion. The total logistics 
costs were \$910 billion, which was equivalent to 8.7\% of the U.S. gross 
domestic product in 2002 (\citet{MacroSys}). Undeniably, logistics
costs are very important for the U.S. economy as well as for a company. 
That is why logistics has received such attention by researchers and managers.   

Both researchers and managers have realized that if the efficiency of logistics activities increases, 
the related costs will decrease and customer satisfaction will increase. 
In addition, since the contribution of logistics activities to the GDP is
large, improvements in the efficiency of the logistics activities will contribute to the 
growth of the U.S. GDP.  Even a small amount of improvement in distribution systems
can result in a substantial amount of savings. Therefore, logistics 
decisions are very important and require professional decisions and
efficient solution techniques to optimize the best interests.
Designing a distribution network requires decisions regarding 
location, routing and scheduling. Companies must decide
where to locate their facilities, how to serve their customers
on routes and how to use their vehicles and/or 
distribution resources efficiently to optimize the flow on their
distribution networks and to optimize storage costs in their warehouses.  
%Our goal is to design a distribution network and to give mathematical 
%formulation of the model that characterizes the logistics system
%and to develop solution methods in order to optimize the decisions for the system. 
%In section \ref{motiv}, we motivate our study by describing simpler 
%distribution systems. 
Our goal is to understand and model the interaction and inter-relationship among facility
location, routing and scheduling decisions  in designing a distribution network. 




%transportation costs were 63\% of the logistics costs while inventory
%carrying costs and warehousing costs are 33\% of the logistics costs. The 
%total logistics costs were \$910 billion in 2002, which is equivalent to 
%8.7\% of the U.S. gross domestic product in 2002 (\citet{}). That us why
%logistics has a lot of attention by researchers and managers.   

%paradan sonra therefore substantial savings can be achieved by
%improving distribution systems even by only a small amount
%GDP informationindan sonra da for US economies is important. 

%Every system has a simple or complicated, physical or 
%non-physical distribution network for which decisions need to be made 
%regarding location, routing and scheduling.  
%
%problems very often and try to develop 
%efficient solution techniques to optimize their best interests. 

%The questions to be answered and decisions to be made.




\section{Outline of the Dissertation}

\begin{comment}
The remainder of the proposal is organized as follows. In chapter \ref{literature}, 
we describe related work in the literature. 
In chapter \ref{problem}, we present edge-based and set partitioning-based formulations 
of the LRS problem and, using the Dantzig Wolfe Decomposition method, we show that they are 
equivalent. We introduce some valid inequalities for the set partitioning model and describe  
our solution approach, a branch-and-price algorithm. In chapter \ref{computationalexp}, 
we present our computational results on the test problems that we have generated. 
%In chapter \ref{proposedRes}, we present our plan for future research. 

**** Add to this part later *******
\end{comment}

\section{Overview of Contributions}

\begin{comment} %First 
In this proposal, we present two new formulations of the LRS problem. Applying the Dantzig Wolfe 
Decomposition method, we show that these two formulations are equivalent. We present two valid 
inequalities and show that the CG rank of these inequalities are 1. Also, we show empirically
that the valid inequalities improve the LP relaxation bound. We develop a branch-and-price
algorithm in which the pricing problem is modeled as an elementary shortest path problem with
resource constraints. We apply three branching rules to optimize the branch-and-price algorithm. 
To date, we can solve instances up to 40 customers and 5 facilities to optimality.   


**** Add to this part later *******

\end{comment}



\begin{comment}
Based on our current study, we propose to improve our branch-and-price algorithm to be able to
solve larger problems. Since our LRS problem is a generalization of some other problems
described in the literature such as MDVRP, LRP and VRPMT, we plan to modify our model and
solution algorithm for the LRS problem to get models and solution algorithms for 
special cases of the LRS problem. 
We propose to find bounds or optimal solutions for MDVRP, LRP and VRPTM models. We propose to 
extend our LRS problem  to an LRS problem with stochastic demands. We will investigate the 
solution methods of the stochastic model. Depending on the remaining time in our study, we may 
consider additional extensions, such as an LRS problem with non-homogeneous fleet and a vehicle 
routing problem with full truckloads.  
\end{comment}
